% ======================================================================
%  TIER 2 KARTI: Reasoning Models
% ======================================================================

\card{Reasoning Models}{düşünme yerleşik modeller}
\cl{Problem} Bazı problemler o kadar çok adımlıdır ki prompt'la Chain-of-Thought
eklemek bile yetmez veya kırılgan kalır --- model yine de acele eder.
\cl{Tanım} \term{Reasoning model}, cevaptan önce kendi içinde uzun bir
``düşünme'' adımı üreten ve bu yeteneği eğitimle kazanmış modeldir. CoT'yi sen
istemeden, içeride yapar.
\cl{Model} Sıradan model hızlı sezgiyle cevaplar; reasoning model ``biraz
düşüneyim'' deyip arka planda müsvedde yapar, sonra konuşur.
\cl{Ne zaman} Karmaşık matematik / mantık / planlama / kod; doğruluğun gecikme ve
maliyetten \emph{önemli} olduğu yerde.
\cl{Tuzak} Basit/hızlı işlere reasoning model koymak (yavaş + pahalı, gereksiz);
ham düşünme zincirini kullanıcıya göstermek.
\cl{İlişkili} \term{Reasoning} (CoT · Reflection · Planning) · \term{Maliyet \& Gecikme} · \term{Hangi Modeli Seçmeli?} \,{\scriptsize\color{softgray}(İlke \#10)}
